\section{Notation}
\label{sec:notation}

\paragraph{General.}
We identify $L$-bit keys with the integers $\{0, 1, \dots, 2^L - 1\}$ via
the standard binary representation. Lexicographic comparison of bit strings
of the same length coincides with numerical comparison of the corresponding
integers, so we use the integer view throughout (this is the natural
formulation when the construction performs arithmetic such as rotations
or reductions modulo a power of two).

\paragraph{Sets and parameters.}
\begin{description}[style=nextline, leftmargin=4cm, labelwidth=3.4cm]
  \item[$L$] Key length in bits.
  \item[{$U = [0, 2^L)$}] Universe of $L$-bit keys.
  \item[$n$] Number of input keys, $n = |S|$.
  \item[$S \subseteq U$] The set of keys in a Range Filter, $|S| = n$.
  \item[$Y = U \setminus S$] Non-keys (all points not in $S$).
  \item[$\mathcal{L}$] Maximum query range length.
  \item[$\varepsilon$] Target false positive rate.
\end{description}

\paragraph{Hash function and fingerprints.}
\begin{description}[style=nextline, leftmargin=4cm, labelwidth=3.4cm]
  \item[$h$] Hash function $h\colon U \to U'$.
  \item[$K$] Fingerprint length in bits, $K = \lceil\log_2(n\mathcal{L}/\varepsilon)\rceil$.
  \item[{$U' = [0, 2^K)$}] Mapped universe.
  \item[{$r = 2^K$}] Mapped universe size.
  \item[$S' = h(S)$] Fingerprints of the input keys.
  \item[$Y' = h(Y)$] Images of non-keys under $h$.
\end{description}

\paragraph{Metrics.}
\begin{description}[style=nextline, leftmargin=4cm, labelwidth=3.4cm]
  \item[FPR] False Positive Rate.
  \item[BPK] Bits Per Key: total filter size divided by $n$.
\end{description}

\paragraph{Structures.}
\begin{description}[style=nextline, leftmargin=4cm, labelwidth=3.4cm]
  \item[ERE] Exact Range Emptiness --- zero-error range filter.
  \item[ARE] Approximate Range Emptiness --- range filter that may have false positives with probability $\leq \varepsilon$.
\end{description}
